\chapter{Introducción}\label{cap:introduccion}

\section{Contexto}

La desinformación se ha consolidado como uno de los principales riesgos
sociotécnicos de la última década, con consecuencias documentadas en
salud pública, procesos electorales y reputación corporativa
\cite{lazer2018science,vosoughi2018spread}. Las grandes plataformas y
los medios de comunicación han desplegado equipos de
\emph{fact-checking}, pero el volumen y la velocidad de difusión de
contenidos hace inviable una revisión manual completa, lo que ha
impulsado la aparición de soluciones asistidas por inteligencia
artificial.

Este Trabajo Fin de Grado en Ingeniería Informática presenta
\textbf{FakeNews Insight}: una plataforma SaaS que permite a un usuario
final analizar un artículo o afirmación textual y obtener (i) un
veredicto clasificatorio basado en un modelo entrenado y (ii) un
informe de verificación a nivel de \emph{claim} con evidencias
externas, generado por un agente de verificación basado en el dataset
FEVER \cite{thorne2018fever}.

\section{Motivación}

El proyecto surge de la combinación de tres motivaciones:

\begin{enumerate}[leftmargin=*,itemsep=0.2em]
  \item \textbf{Investigadora.} La parte estadística y de modelado se
        desarrolla en una memoria independiente del Grado en
        Estadística del mismo autor, lo que permite enfocar este TFG
        de Informática en los aspectos puramente \emph{software} y de
        sistemas.
  \item \textbf{Producto.} Se busca diseñar e implementar una
        plataforma usable, multicliente, multi-idioma y con planes de
        suscripción gestionados, no un mero \emph{prototipo}
        académico.
  \item \textbf{Industrial.} Aplicar metodologías de ingeniería de
        software actuales: arquitectura desacoplada, despliegue
        contenedorizado, control de acceso basado en filas, autenticación
        federada, observabilidad básica y \emph{testing} automatizado.
\end{enumerate}

\section{Objetivos}

\begin{enumerate}[leftmargin=*,itemsep=0.2em]
  \item Diseñar e implementar un \textbf{backend} REST en FastAPI con
        endpoints de inferencia, gestión de cuenta, historial,
        \emph{billing} con Stripe y verificación de afirmaciones.
  \item Desarrollar un \textbf{frontend} SPA en React con
        arquitectura unidireccional UI $\to$ Estado $\to$ Servicios,
        soporte i18n (es/en), tests Vitest y diseño \emph{mobile-first}.
  \item Construir una \textbf{extensión de navegador MV3} que reutilice
        la API y la autenticación de la plataforma.
  \item Integrar un \textbf{agente de verificación} basado en
        recuperación de evidencias y NLI, expuesto como servicio
        diferenciado del plan \emph{Super Pro}.
  \item Implantar un modelo de datos sobre Supabase con \textbf{Row
        Level Security} y migraciones SQL versionadas.
  \item Empaquetar todo en \textbf{Docker Compose} y dejar
        configuración lista para despliegue en una VPS.
\end{enumerate}

\section{Estructura de la memoria}

El capítulo~\ref{cap:estado-arte} revisa soluciones similares y
tecnologías candidatas. El capítulo~\ref{cap:requisitos} formaliza
requisitos funcionales y no funcionales. El capítulo~\ref{cap:arquitectura}
expone la arquitectura general. Los
capítulos~\ref{cap:backend},~\ref{cap:frontend} y~\ref{cap:extension}
documentan los tres componentes principales. El
capítulo~\ref{cap:agente} describe el agente de verificación. El
capítulo~\ref{cap:billing} cubre el subsistema de \emph{billing} y
planes. El capítulo~\ref{cap:seguridad} desarrolla seguridad y RLS.
El capítulo~\ref{cap:pruebas} resume la estrategia de pruebas. El
capítulo~\ref{cap:despliegue} describe el despliegue. Los
capítulos~\ref{cap:resultados} y~\ref{cap:conclusiones} cierran con
resultados y conclusiones.
